\section{Human Evaluation Study}
\label{sec:human-study}

Both experiments so far rely on language models to score the output of a
language model. The third study asks people.

This study cannot validate the 35-pull-request ranking: the stimuli are a
different sample, the graph block's edges are resolved from import statements rather than
from the Joern CPG, the prompt is stricter, and the criteria are reworded for non-expert readers. What it
tests is whether human readers perceive the specific difference the graph
block produces.

\subsection{Stimuli and Design}

Six pull requests were drawn from three Python libraries, with two from each:
\texttt{psf/requests}, \texttt{pallets/flask}, and \texttt{pallets/click}.
The restriction to one language and to well-known libraries reduces the prior
repository knowledge required of raters.

Each pull request is presented with two reviews: \baseline{} (strict prompt,
no graph block) and \kg{} (a matched prompt that differs only in its
references to the available context, plus the knowledge-graph block). The on-screen diff is shared reference material for both
arms. Reviews are shown as A and B, and participants are not told which mode
produced either review. The side on which each mode appears is randomised per
rater.

The design is fully crossed: every rater sees all six pull requests. Each
task presents a plain-language summary, the two reviews side by side, and
the diff, which is expanded by default. The assisted interface also provides
arm-neutral context, review summaries, file-existence tooltips that do not
certify affectedness, and optional difference highlighting that is off by
default. The rater makes six per-criterion choices (A, B, both, or neither),
states an overall preference (A, equal, or B), optionally explains the choice
in free text, and rates the task's difficulty.
Table~\ref{tab:study-criteria} lists the criteria.

\begin{table}[htbp]
    \centering
    \caption{The six comparison criteria as worded in the instrument. H1 is
             the primary endpoint. IDs are prefixed H to distinguish them
             from the 25-criterion rubric used in Experiments~1 and~2; the
             closest rubric counterparts are F3/F4 for H1, F2 for H2, T3 for
             H3, and R1 for H5.}
    \label{tab:study-criteria}
    \small
    \begin{tabular}{llp{7.4cm}}
        \toprule
        ID & Category & Question put to the rater \\
        \midrule
        H1 & Functionality  & Which review better names the specific functions, classes, or files that this change could affect? \\
        H2 & Functionality  & Which review better describes specific situations where the changed code could fail or misbehave? \\
        H3 & Tests          & Which review better points to specific test files, or says which tests should be added or updated? \\
        H4 & Quality        & Which review better explains \emph{why} something is a problem, rather than just saying it is one? \\
        H5 & Readability    & Which review better comments on how clear and well-organised the code is? \\
        H6 & Completeness   & Which review better covers the important issues in this PR, without obvious gaps? \\
        \bottomrule
    \end{tabular}
\end{table}

\subsection{Analysis Plan}

Choices are scored directionally: a pick of the \kg{} side counts one, a pick
of the \baseline{} side counts zero, and ``both'', ``neither'', or ``equal''
count one half. ``Both'' and ``neither'' share the same directional score but
remain separate in the raw-choice descriptives. A rater's score for a criterion
is the mean over their six tasks; the unit of analysis is the rater and the
null value is~$0.5$.

The primary endpoint is H1 (naming affected components), tested against~$0.5$
with a two-sided Wilcoxon signed-rank test at $\alpha = 0.05$, reported with
a rank-biserial correlation and a bootstrap confidence interval. This
criterion was chosen because the LLM-judged result localises its effect
there, and it is where the two arms differ mechanically. Overall usefulness
is secondary and explicitly estimation-only: power simulations show it
underpowered at the intended sample size. The remaining criteria are
exploratory and Holm-corrected.

The analysis plan also pre-registered four diagnostic checks concerning
file-count association, aid use, free-text rationales, and task time. They were
not implemented in the final analyser and are not used to support any thesis
claim; this omission is reported as a limitation.

The recruitment target is 18--20 complete raters, with 15 as a lower-bound
checkpoint. Only raters who complete all six tasks enter the confirmatory
analysis. The final export exceeded that upper target.
Section~\ref{sec:results-human} reports all 27 complete raters. The
interpretation rules were fixed before data collection: a positive result
licenses the claim that raters perceive the graph arm's better naming of
affected code; a
null result is reported as a boundary finding and the plan forbids re-running
with new stimuli.

\subsection{Ethics}

Participation begins with informed consent. The instrument records a
self-chosen identifier, optional free text, role and experience descriptors,
criterion and overall choices, difficulty, task time, GitHub-link clicks, and
whether highlighting was enabled; it collects neither email nor IP address.
Processing rests on consent under Art.~6(1)(a) GDPR. Data are retained until
examination is complete and deleted no later than 31~December~2027. Only
aggregated or de-identified results appear in this thesis. No formal ethics
board review was required under the university's guidelines for this category
of research; the study collects no personally identifying information and
presents no risk beyond standard professional activity.
